\chapter{Metric Independence and Cohomological Observables}
\label{ch:tqft-metric-independence}

\ConventionsEuclidean

A topological QFT assigns quantities to manifolds and bordisms in a way that
is insensitive to continuous changes of geometric representatives after the
allowed background structures are fixed.  A cohomological QFT realizes this
insensitivity by an odd differential \(Q\) acting on fields, states, and
operators.

\section{Metric Variation}

Let \(M\) be a closed oriented \(D\)-manifold and let \(g\) be a Riemannian
metric.  Suppose the theory has correlation functions
\(\langle\mathcal O_1\cdots\mathcal O_n\rangle_g\) for metric-dependent
representatives of local or extended observables.  The stress tensor is
defined by metric variation of the generating functional:
\[
  \delta_g\langle X\rangle_g
  =
  -\frac12
  \int_M\dd^Dx\,\sqrt g\,\delta g^{\mu\nu}(x)\,
  \langle T_{\mu\nu}(x)X\rangle_g
  +\hbox{contact terms from metric-dependent insertions}.
\]
This formula is a definition of the response to background geometry.  The
contact terms are part of the operator definition and must be retained when
insertions collide with the support of \(\delta g\).

\begin{definition}[Cohomological topological sector]
\label{def:cohomological-topological-sector}
A cohomological topological sector consists of a graded vector space of
states or observables, an odd differential \(Q\) with \(Q^2=0\), a class of
\(Q\)-closed observables, and tensors \(G_{\mu\nu}\) such that
\[
  T_{\mu\nu}=\{Q,G_{\mu\nu}\}
\]
inside correlation functions, including the contact-term prescription for
metric variation.
\end{definition}

\begin{proposition}[Metric independence in \(Q\)-cohomology]
\label{prop:q-cohomology-metric-independence}
Assume the hypotheses of Definition~\ref{def:cohomological-topological-sector}
and suppose \(X=\mathcal O_1\cdots\mathcal O_n\) is a product of
\(Q\)-closed observables whose metric-dependence is itself \(Q\)-exact.  If
the functional expectation satisfies the Ward identity
\(\langle \{Q,Y\}\rangle_g=0\) for all insertions \(Y\) produced by the metric
variation, then \(\delta_g\langle X\rangle_g=0\).
\end{proposition}

\begin{proof}
Substituting \(T_{\mu\nu}=\{Q,G_{\mu\nu}\}\) into the metric variation gives
an integral of \(\langle\{Q,G_{\mu\nu}\}X\rangle_g\), plus the declared
metric-variation contact terms.  Since \(QX=0\) and \(Q\) is an odd
derivation, this is \(\langle\{Q,G_{\mu\nu}X\}\rangle_g\) up to the signs
fixed by the grading.  The contact terms are \(Q\)-exact by hypothesis.  The
Ward identity therefore makes every term vanish.
\end{proof}

\section{Functorial Target}

The Atiyah--Segal form of a topological field theory assigns a vector space
\(Z(\Sigma)\) to a closed \((D-1)\)-manifold \(\Sigma\) and a linear map
\[
  Z(M):Z(\Sigma_{\rm in})\to Z(\Sigma_{\rm out})
\]
to a bordism \(M:\Sigma_{\rm in}\to\Sigma_{\rm out}\), with compatibility
under gluing:
\[
  Z(M_2\circ M_1)=Z(M_2)Z(M_1).
\]
Extended versions refine this assignment to manifolds with corners and
higher-categorical targets.  The cohomological construction above is one
source of such functors after gauge fixing, anomaly cancellation, boundary
conditions, and finite-dimensionality hypotheses have been supplied.

\begin{openproblem}[From cohomological path integrals to extended functors]
\label{op:cohomological-to-extended-functor}
Given a cohomological gauge theory with \(T_{\mu\nu}=QG_{\mu\nu}\), determine
the exact hypotheses under which its correlation functions define an
extended topological functor: compactness or properness of the moduli
problem, anomaly cancellation, boundary conditions, gluing analysis,
orientation data, and compatibility with local operator insertions.
\end{openproblem}
